\section{Introduction}

Consider a long-running agent task that crosses two sessions. The first
session changes a repository, starts remote work, records one failed check, and
ends before the task is complete. A later participant opens the local product
and asks what happened and what action is safe now. The local interface, remote
runtime, repository, build system, and retained session evidence may each
contain valid facts. They need not contain the same facts, reach the same cut,
or order concurrent facts in the same way.

The KFD Foundation Triad describes how such work can remain one shared world
across participants and time \cite{kfdFoundation2026}:

\begin{equation}
\text{Facts} \rightarrow \text{Trust} \rightarrow \text{Cooperation}
\rightarrow \text{New Facts}.
\end{equation}

KFD-1 preserves stable reference, KFD-2 binds reliance to evidence, and KFD-3
makes relevant value and constraints visible to a participant before action.
Cooperation then creates consequences that return through review into a
successor fact world. That loop solves a continuity problem, but it creates a
new pressure: no participant consumes the entire world from every position.

The pressure appears separately at each layer of the triad. If a selected or
transformed fact cut is presented without perspective, selection can masquerade
as reality and violate the stable reference required by KFD-1. If a claim does
not identify which cut and projection support it, a later participant cannot
reconstruct the warrant required by KFD-2. If a view hides whose consequence
position, available actions, and known gaps it represents, a participant cannot
inspect the value and constraints needed by KFD-3; cooperation degrades into
guessing, obedience, or repeated reconstruction.

This yields the first derived pressure from the Foundation Triad:

\begin{equation}
\text{non-drifting facts} + \text{multiple participants}
\Longrightarrow \text{declared perspective}.
\end{equation}

KFD-4 names the corresponding rule: timelines must declare their observer
\cite{kfd2026observer}. An observer-declared timeline does not deny facts or
causality. It states which observer position is represented, which facts were
accepted, how the view was projected, which causal relations dominate that
projection, what consequences or actions the view is meant to support, and
where evidence remains incomplete. The result is not a view from nowhere. It
is a reproducible claim about how one view was formed from a shared fact world.

The core rule leads to a second question only after the perspective is defined:
what happens when the observer implementation itself restarts, upgrades, or is
replaced? A stable product label does not prove that a new process accepted the
same facts or applied the same projection semantics. We therefore distinguish a
logical observer from its incarnations and treat continuity between them as an
evidence-bearing transition. This is a derived refinement of the core rule, not
a prerequisite for first understanding why perspective must be declared.

This paper makes five contributions:

\begin{itemize}[leftmargin=*]
  \item a derivation of observer declaration from the fact, trust, and
        cooperation responsibilities of the KFD Foundation Triad;
  \item a minimal model over accepted fact cuts, causal constraints, projection
        policy, consequence position, degradation, and verification;
  \item an incarnation and continuity-witness refinement for observers that
        fail, restart, or change implementation;
  \item a mapping from that generic model to Kungfu's Episode fact authority,
        exact-cut readiness, fenced recovery, and product surfaces; and
  \item bounded implementation evidence together with an explicit account of
        the observer-level evaluation that remains incomplete.
\end{itemize}

Section~\ref{sec:background} develops the problem through the running example
and distributed-systems context. Section~\ref{sec:model} defines the generic
model without requiring a Kungfu-specific causal object. Section~\ref{sec:design}
then introduces Episodes and runtime mechanisms as one concrete realization.
Section~\ref{sec:evaluation} separates evidence for that substrate from the
still-open end-to-end observer claim, and Section~\ref{sec:related} compares
related work.
